\documentclass[11pt]{article}
\usepackage[margin=1in]{geometry}
\usepackage{amsmath,amssymb,amsthm,mathtools}
\usepackage[T1]{fontenc}
\usepackage[expansion=false]{microtype}
\usepackage[hidelinks]{hyperref}
\newtheorem{theorem}{Theorem}[section]
\newtheorem{lemma}[theorem]{Lemma}
\newcommand{\ddc}{\mathrm d\mathrm d^c}
\newcommand{\NP}[1]{\left\langle#1\right\rangle}
\title{Divisorial smoothing and approximation-dependent adjoint probabilities}
\author{}
\date{}
\begin{document}
\maketitle
\begin{abstract}
Smooth positive approximations of a logarithmic potential need not recover
its relative non-pluripolar adjoint. On $\mathbb P^1\times\mathbb P^1$ we
construct an admissible ordinary model reference and a singular solution
whose supporting probability is unique. Two monotone smooth regularizations
of this same solution retain a fixed positive lower bound, converge in
capacity and smoothly off the divisor, and have logarithmic densities
converging in every finite $L^p$. Their normalized classical adjoint
probabilities nevertheless have different weak limits, both charging the
divisor. We compute the limits and the residual weak linearization explicitly.
The obstruction concerns unrestricted smoothing of poles; it does not apply
to approximation classes preserving singularity type or relative mass.
\end{abstract}

\section{The assertion and the counterexample}

Use $\ddc=(i/(2\pi))\partial\bar\partial$ and exterior powers without
divided factorials. Let $X=\mathbb P^1\times\mathbb P^1$. Denote the
pulled-back Fubini--Study forms, each of integral one on its factor, by
$\omega_1,\omega_2$. Set
\[
 \theta=2\omega_1+2\omega_2,\qquad
 \mathrm dV=\omega_1\wedge\omega_2,\qquad
 \beta=\omega_1+2\omega_2.
\]
Fix $p=0$ in the affine coordinate $z$ on the first factor. The globally
defined squared norm of a section of $\mathcal O(1)$ vanishing at $p$ is
\[
 r=\frac{|z|^2}{1+|z|^2},\qquad
 D=\{p\}\times\mathbb P^1,\qquad \psi=\log r.
\]
The canonical value of $\psi$ is $-\infty$ on $D$.

We use ordinary bounded equivalence of singularities: $v\sim\psi$ means
that $v-\psi$ is bounded off $D$, with their common pole values retained.
The envelope is the upper regularization of genuine $\theta$-psh
minorants, and
\[
 P_\theta[\psi]=
 \left(\lim_{C\to\infty}P_\theta(\min\{\psi+C,0\})\right)^*.
\]
For an admissible model reference $\phi$, write
$B_\phi(F)=\log c_\phi(F)$, where the relative equation is
\[
 \operatorname{MA}_{\theta}^{\mathrm{NP}}(v)
      =c_\phi(F)e^F\mathrm dV,\qquad v\sim\phi.
\]
The problem supplies solvability and scalar uniqueness for bounded $F$
in this admissible regime. Its supporting set consists of probabilities
$\nu$ satisfying
\[
 B_\phi(G)\le B_\phi(F)-\int_X(G-F)\,\mathrm d\nu
 \quad\text{for all continuous real }G.
\]

\begin{theorem}\label{thm:obstruction}
The reference $\phi=P_\theta[\psi]$ is an ordinary model, is boundedly
equivalent to $\psi$, has supremum zero, and has a genuine
$P_{(3/4)\theta}(\phi)$. The potential $u=\psi$ is a relative solution
at $F=0$, with scalar $4$ and singleton supporting set $\{\mathrm dV\}$.
Its normalized non-pluripolar adjoint is $\beta/4$.

For every smooth real $g$ on the second factor with
$\ddc g\ge-\omega_2$, define
\[
 u_\delta^g=\log(r+\delta e^g),\qquad
 \Omega_\delta^g=\theta+\ddc u_\delta^g,\qquad \delta>0.
\]
These are globally smooth, decrease canonically to $u$ as
$\delta\downarrow0$, and satisfy $\Omega_\delta^g\ge\omega_1+\omega_2$.
The convergence is in $L^1$, in compact K\"ahler capacity, and in
$C^\infty$ on every compact subset of $X\setminus D$.
Their normalized classical adjoints are
\[
 \rho_\delta^g=\frac18,\qquad
 Q_\delta^g=\frac18\Omega_\delta^g,\qquad
 \nu_\delta^g=\frac18(\Omega_\delta^g)^2.
\]
Their weak limits satisfy
\begin{align}
 Q_\delta^g&\rightharpoonup\frac18(\beta+[D]),\label{eq:currentlimit}\\
 \nu_\delta^g&\rightharpoonup
 \frac12\mathrm dV+\frac12[D]\wedge\omega_2
                         +\frac18[D]\wedge\ddc g.\label{eq:measurelimit}
\end{align}
In particular $g=0$ and a sufficiently small nonconstant $g$ give
different probability limits. Neither limit is the relative supporting
probability $\mathrm dV$.

The smooth logarithmic densities
\[
 F_\delta^g=\log\frac{(\Omega_\delta^g)^2}{4\mathrm dV}
\]
converge to zero in every finite $L^p(\mathrm dV)$, although
$\int_Xe^{F_\delta^g}\mathrm dV=2$ for every $\delta$.
\end{theorem}

Thus canonical, monotone and capacity convergence, even together with
uniform lower positivity and finite-$L^p$ convergence of logarithmic
densities, does not give an approximation-independent relative adjoint
under unrestricted smoothing of a pole. The smooth approximants are
bounded and do not preserve the singular reference type.

\section{An explicit admissible model reference}

Poincar\'e--Lelong gives
\begin{equation}\label{eq:PL}
 \ddc\psi=[D]-\omega_1,\qquad
 \theta+\ddc\psi=\beta+[D],\qquad \beta^2=4\mathrm dV.
\end{equation}
The convention can also be checked directly from
$\ddc\log|z|^2=[0]$ and $\omega_1=\ddc\log(1+|z|^2)$.

Here the model envelope admits an elementary formula. Put
\[
 x=\log|z|^2,\qquad \chi(x)=\log(1+e^x),\qquad
 W(x)=\begin{cases}
         x+2\log2,&x\le0,\\
         2\chi(x),&x\ge0.
       \end{cases}
\]
Then
\begin{equation}\label{eq:model}
 \phi=W-2\chi.
\end{equation}
Indeed $W$ is convex, with matching derivative $1$ at zero and limiting
slopes $1,2$ at the two ends. The radial function $W-2\chi$ is therefore
genuine $2\omega_1$-psh on the first factor, with a logarithmic pole at
$p$ and its canonical value at infinity. It is independent of the
second factor, and hence $\theta$-psh on $X$. Direct subtraction gives
\[
 \psi\le\phi\le0,\qquad 0\le\phi-\psi\le2\log2.
\]

To check the envelope identity, let
$v_C=P_\theta(\min\{\psi+C,0\})$, $C\ge0$. The obstacle and $\theta$
are invariant under rotations on the first factor and the transitive
$\mathrm{SU}(2)$ action on the second. Pullback preserves the complete
class of competitors, so $v_C$ has the same invariances. It is therefore
radial on the first factor and independent of the second. On the affine
punctured chart, $V_C=2\chi+v_C$ is convex. The inequality
$v_C\le\psi+C$ gives
$V_C(x)\le\chi(x)+x+C$, so its limiting slope at $-\infty$ is at
least one. Every subsequent slope is at least one. For $x<0$ this yields
\[
 V_C(x)\le V_C(0)+x\le2\log2+x.
\]
For $x\ge0$, the obstacle gives $V_C\le2\chi$. Thus $v_C\le\phi$.
When $C\ge2\log2$, the displayed bounds make $\phi$ a competitor,
so $v_C=\phi$. This proves $P_\theta[\psi]=\phi$.

The operator $P_\theta[\cdot]$ is unchanged by a bounded change of
its input: a two-sided bound on the inputs squeezes their shifted
obstacle envelopes between those with shifted parameters $C$.
Since $\phi\sim\psi$, it follows that $P_\theta[\phi]=\phi$.
Its supremum is zero. Finally
\[
 \tfrac34\theta+\ddc\psi
       =\tfrac12\omega_1+\tfrac32\omega_2+[D]\ge0,
 \qquad \psi\le\phi.
\]
Thus $\psi$ is a genuine minorant for $P_{(3/4)\theta}(\phi)$;
the latter is bounded above by zero and is not identically $-\infty$.
The required admissibility condition is verified.

\section{Residual products, the scalar and the singular adjoint}

If $v\sim\phi$, then $w=v-\psi$ is bounded on $X\setminus D$ and
$\beta+\ddc w\ge0$ there. Add a local smooth potential for the closed
form $\beta$. The resulting bounded psh function extends canonically
across $D$ by the removable-singularity theorem
\cite[Chapter I, Theorem 5.24]{Dem}. The local extensions agree off
$D$, hence everywhere as canonical representatives. Thus $w$ is a
global bounded $\beta$-psh function. Conversely, such a $w$ gives the
genuine potential $\psi+w$ of the prescribed ordinary type, by
\eqref{eq:PL}. Write $T=\beta+\ddc w$.

All products in this paragraph are bounded Bedford--Taylor products.
Their recursive definition, decreasing continuity and integration by
parts are those of \cite[Section 2]{BT82} and
\cite[Chapter III, Section 3]{Dem}. Global integration by parts is
\cite[Theorem 1.14]{BEGZ}, applied with the empty exceptional set,
bounded quasi-psh functions and the closed positive current $T+\beta$.
We verify directly the particular
noncharging and mass statements needed here. Shift $w$ so that
$0\le w\le M$, and set
\[
 f_t=\max\{\psi/t,-1\},\qquad t>0.
\]
This is a bounded continuous quasi-psh cutoff with
$\ddc f_t\ge-\omega_1/t$. It equals $-1$ on $D$ and tends to zero
off $D$. The bounded-current identity
$T^2-\beta^2=\ddc(w(T+\beta))$ gives
\begin{align*}
 \int_X(-f_t)T^2
 &=\int_X(-f_t)\beta^2
             +\int_Xw\,\ddc(-f_t)\wedge(T+\beta)\\
 &\le\int_X(-f_t)\beta^2
             +\frac Mt\int_X\omega_1\wedge(T+\beta).
\end{align*}
The final mixed mass is $4$ by Stokes. Dominated convergence for the
smooth first term shows $T^2(D)=0$. Stokes in the same current identity
also gives $\int_XT^2=\int_X\beta^2=4$. Non-pluripolar locality off
$D$, followed by zero extension across $D$, therefore gives
\[
 \operatorname{MA}_{\theta}^{\mathrm{NP}}(v)=T^2,
 \qquad \int_X\operatorname{MA}_{\theta}^{\mathrm{NP}}(v)=4.
\]
Integrating the relative equation and using the supplied solvability
and scalar uniqueness yields, for every bounded real $F$,
\[
 c_\phi(F)=\frac4{Z_F},\qquad
 B_\phi(F)=\log4-\log Z_F,\qquad Z_F=\int_Xe^F\mathrm dV.
\]
For continuous $F$, Jensen's inequality puts
$\nu_F=e^F\mathrm dV/Z_F$ in the supporting set. If $\nu$ is any
other support, apply its inequalities to $F\pm th$ and differentiate
at zero for every continuous $h$. This gives
$\int h\,\mathrm d\nu=\int h\,\mathrm d\nu_F$, hence
\[
 \mathcal A_\phi(F)=\{e^F\mathrm dV/Z_F\}.
\]

At $u=\psi$, the non-pluripolar first and second products are $\beta$
and $\beta^2$. The ordinary smooth multiplication
\[
 Q=\rho\NP{\theta+\ddc u}=\tfrac14\beta,\qquad \rho=\tfrac14
\]
satisfies $\ddc Q=0$ and
$\int_X\rho\NP{(\theta+\ddc u)^2}=1$. Its associated probability
is $\rho\beta^2=\mathrm dV$, the unique support at $F=0$.
For every smooth $h$, the derivative of the non-pluripolar square
along $u+th$ is the meaningful measure-valued operator
\[
 \mathcal L_uh=2\ddc h\wedge\beta.
\]
For small $t$ the residual form remains positive, and Stokes gives
$\int_X\rho\mathcal L_uh=0$. No product with an undefined value
of $u$ on $D$ is taken.

\section{Positive regularizations and the normal concentration}

Set $\alpha=\log r$, $\gamma=\ddc g$, and, on $X\setminus D$,
\[
 \lambda=\frac r{r+\delta e^g},\qquad
 h=\alpha-g-\log\delta,\qquad
 \mathrm d^c=\frac{i}{4\pi}(\bar\partial-\partial).
\]
The exact chain rule gives
\begin{align*}
 \Omega_\delta^g&=A_\delta+R_\delta,\\
 A_\delta&=(2-\lambda)\omega_1+2\omega_2+(1-\lambda)\gamma,\\
 R_\delta&=\lambda(1-\lambda)\mathrm dh\wedge\mathrm d^ch.
\end{align*}
Here $R_\delta$ is positive of rank one and
$A_\delta\ge\omega_1+\omega_2$. Smoothness extends this positivity
across $D$. Each $\Omega_\delta^g$ is closed and
\[
 \int_X(\Omega_\delta^g)^2=\int_X\theta^2=8.
\]
Its smooth adjoint equation is
$\ddc\rho\wedge\Omega_\delta^g=0$. The maximum principle on the
connected compact manifold makes $\rho$ constant, and normalization
therefore gives $\rho=1/8$.

Since $R_\delta^2=0$,
\[
 (\Omega_\delta^g)^2=A_\delta^2+2A_\delta\wedge R_\delta.
\]
The first term converges to $\beta^2=4\mathrm dV$ by dominated
convergence. Put
$Q_N=\mathrm d\alpha\wedge\mathrm d^c\alpha$ and
$Q_T=\mathrm dg\wedge\mathrm d^cg$. The mixed normal--tangential
terms vanish when wedged with the diagonal form $A_\delta$. Thus
\begin{align*}
 A_\delta\wedge R_\delta
 &=\lambda(1-\lambda)Q_N\wedge
                  \bigl(2\omega_2+(1-\lambda)\gamma\bigr)\\
 &\quad+\lambda(1-\lambda)(2-\lambda)\omega_1\wedge Q_T.
\end{align*}
The second line tends to zero by bounded dominated convergence.
The two normal concentration identities are
\begin{equation}\label{eq:spikes}
 \lambda(1-\lambda)Q_N\rightharpoonup[D],\qquad
 \lambda(1-\lambda)^2Q_N\rightharpoonup\tfrac12[D].
\end{equation}

To see the constants directly, fix a second-factor point, write
$\delta'=\delta e^g$, and rescale $z=\sqrt{\delta'}\,\zeta$.
The first normal measure converges to
\[
 \frac{i}{2\pi}\frac{\mathrm d\zeta\wedge\mathrm d\bar\zeta}
                              {(1+|\zeta|^2)^2}.
\]
Its total mass is one. In the radial variable $y=|\zeta|^2$ its measure
is $\mathrm dy/(1+y)^2$, and $\lambda\to y/(1+y)$. Consequently
$\lambda$ is uniform on $[0,1]$ for this limiting spike, and
the mean of $1-\lambda$ is $1/2$. On a fixed small first-factor chart,
the rescaled tails are bounded by a constant times the tail of
$\mathrm dy/(1+y)^2$; outside the chart the mass is $O(\delta')$.
These bounds justify arbitrary continuous test functions and are
uniform on the second factor, since $\delta'/\delta$ is bounded
above and below. This proves \eqref{eq:spikes} as global currents.

Combining the identities yields
\[
 (\Omega_\delta^g)^2\rightharpoonup
       4\mathrm dV+4[D]\wedge\omega_2+[D]\wedge\gamma,
\]
which gives \eqref{eq:measurelimit}. The limit is positive because
$4\omega_2+\gamma\ge3\omega_2$, and has total mass one because
$\int_{\mathbb P^1}\gamma=0$. Its mass on $D$ is $1/2$.
A sufficiently small nonconstant $g$ has $\gamma\ne0$: a globally
harmonic real function on $\mathbb P^1$ is constant. Hence its limit
differs from the one with $g=0$.

\section{Convergence and the boundary of the obstruction}

The explicit formula gives pointwise monotone convergence with the
canonical pole values and smooth convergence off $D$. Moreover
\[
 0\le u_\delta^g-\psi
    \le\log\left(1+\frac{\delta e^{\max g}}r\right)
 \quad\text{on }X\setminus D.
\]
Logarithmic integrability near the smooth divisor proves $L^1$
convergence. Distributional differentiation gives
\eqref{eq:currentlimit}.

For completeness, capacity convergence follows from the same bounded
cutoffs used above. Put $\omega=\omega_1+\omega_2$ and define
$\operatorname{Cap}_\omega(E)$ as the supremum of
$\int_E(\omega+\ddc v)^2$ over $\omega$-psh $-1\le v\le0$.
With $T_v=\omega+\ddc v$, the cutoff calculation gives uniformly in $v$
\[
 \int_X(-f_t)T_v^2
     \le\int_X(-f_t)\omega^2+\frac2t.
\]
It follows that $\operatorname{Cap}_\omega(\{\psi<-t\})\to0$.
For every $\eta>0$,
\[
 \{u_\delta^g-\psi>\eta\}\subset
   \left\{r<\frac{\delta e^{\max g}}{e^\eta-1}\right\},
\]
with the polar set ignored in capacity. This proves the asserted
convergence, including that $D$ itself has capacity zero.

Finally $\Omega_\delta^g\ge\omega$ implies
$F_\delta^g\ge-\log2$. On $X\setminus D$, the bounded derivatives
of $g$ and the inequality $|\mathrm d\log r|_\omega^2\le C/r$ give
\[
 \Omega_\delta^g\le C(1+1/r)\omega,\qquad
 |F_\delta^g|\le C+2|\log r|.
\]
The majorant is in every finite $L^p(\mathrm dV)$, while
$F_\delta^g\to0$ off $D$. Dominated convergence proves the stated
logarithmic-density convergence. Yet the smooth mass identity gives
\[
 \int_Xe^{F_\delta^g}\mathrm dV=2\ne1=\int_Xe^0\mathrm dV.
\]
The exponentiated densities are therefore not uniformly integrable.

The limit in \eqref{eq:currentlimit} cannot be $\rho\beta$ for a
locally integrable function $\rho$: multiplication by the smooth
form $\beta$ gives locally integrable coefficients, whereas $[D]$
has a nonzero divisorial coefficient measure. Thus even the common
limit of the weighted currents does not recover the normalized
non-pluripolar cofactor $\beta/4$; its associated top products retain
additional information depending on the smoothing path.

This disproves the assertion that unrestricted smooth positive
regularizations determine an approximation-independent relative
adjoint compatible with a singleton support. It does not disprove
an extension for approximations preserving the pole type, relative
mass or appropriate uniform integrability of the exponentiated
densities. No uniform $L^\infty$ bound on the logarithmic densities,
fixed-type approximation theorem or general singular weak-adjoint
existence theorem is claimed.

\begin{thebibliography}{3}
\bibitem{BT82} E.~Bedford and B.~A.~Taylor,
\emph{A new capacity for plurisubharmonic functions},
Acta Mathematica 149 (1982), 1--40,
\url{https://doi.org/10.1007/BF02392348}.
\bibitem{BEGZ} S.~Boucksom, P.~Eyssidieux, V.~Guedj and A.~Zeriahi,
\emph{Monge--Amp\`ere equations in big cohomology classes},
Acta Mathematica 205 (2010), 199--262, Theorem 1.14,
\url{https://www.math.univ-toulouse.fr/~guedj/fichierspdf/BEGZ10.pdf}.
\bibitem{Dem} J.-P.~Demailly,
\emph{Complex Analytic and Differential Geometry},
author-hosted book, Chapters I and III,
\url{https://people.math.harvard.edu/~siu/2dmml.pdf}.
\end{thebibliography}
\end{document}
